\originalsection{边界受力与能量的计算步骤}
在表面场跃变处, “把电荷乘上电场”先要回答电场是否包含这片电荷自身的场. 取局部近似平面的薄带电层, 两侧的自身场为 $\pm\sigma\nn/(2\eps)$; 来自其余电荷的场在层的两侧相同. 将两侧总场平均, 自身场抵消, 得 $\ee_\mathrm{other}=(\ee_\mathrm{out}+\ee_\mathrm{in})/2$. 因而面力密度为 $\sigma\ee_\mathrm{other}$. 对导体 $\ee_\mathrm{in}=0$, 真空外场 $\ee_\mathrm{out}=\sigma\nn/\eps$, 压力才是 $\sigma^2/(2\eps)$. 这个论证使用薄层极限和局部平面近似; 尖角处应先求可积的总力, 不强行给有限点值压力.

场能中的 $1/2$ 也应从建立过程判断. 把最终密度 $\rho$ 缓慢从零增至 $\lambda\rho$, $0\le\lambda\le1$. 线性真空理论使已有电势为 $\lambda V$, 新添电荷为 $\rho\dd\lambda$. 因而
\[
 U=\int_0^1\dd\lambda\int\lambda V\rho\dd^3x
 =\frac12\int\rho V\dd^3x.
\]
用 $\rho=-\eps\nabla^2V$ 后, 格林第一恒等式给
$\int V\nabla^2V=\oint V\partial_nV-\int|\nabla V|^2$.
对局域有限能量分布, 无穷远项消失. 若求的是一个有限区域而边界有非零势, 该边界项一般不能删. 例如电容器场能可以只在有场区域积分, 但 $\frac12\int\rho V$ 的 $\rho$ 要包含导体面的电荷.

固定电压的虚功可逐步算出. $Q=CV$ 给 $\dd Q=V\dd C$, 场能变化为 $\dd U_\mathrm{field}=V^2\dd C/2$, 电源损失为 $\dd U_\mathrm{source}=-V\dd Q=-V^2\dd C$. 总势能下降 $-V^2\dd C/2$, 故沿使 $C$ 增加的方向有正机械力. 固定电荷时没有电源交换, $\dd U=-Q^2\dd C/(2C^2)$, 所得方向相同. 本科作业2题4和作业6题6分别用这两个结论; 读者应先判定保持的是 $Q$ 还是 $V$, 再微分.

\originalsection{法向边界的相容条件}
若给定 $\partial_nV=h$ 而非边界势, 对泊松方程积分得到必要条件
$\oint_{\partial\Omega}h\dd a=-\eps^{-1}\int_\Omega\rho\dd^3x$. 它说指定向外的通量不能与总电荷矛盾. 两个同源、同法向数据解之差 $u$ 有 $\partial_nu=0$, 同一个能量积分仍给 $\nabla u=0$, 但无法再把常数定为零. 指定一个参考点电势或平均值即可除去该自由度. 这解释了为何“只求电场”与“求具有指定零点的势”在边界信息上不同.
